% INSERCIÓN FOCAL AUTORIZADA: no introduce un capítulo ni el concepto de radión.
% Incluir dentro de la sección constitutiva, después de la inversión cúbica
% y de la declaración de las cartas dimensionales. Etiquetas fuera de grupos
% de prefijado automático. Conserva íntegramente el bloque geométrico anterior.
\subsection{Reconstrucción de la forma de acción desde la respuesta constitutiva}
\label{iii:subsec:puente-forma-LC}

La respuesta constitutiva y la realización inductiva--capacitiva reciben
el mismo par angular construido desde APP--TRIT--TPK. La primera evalúa
funciones racionales de los dos canales; la segunda utiliza la razón de
sus exponentes para determinar la forma de la elipse de acción. La
inversión cúbica permite componer ambas lecturas sin identificar sus
cocientes numéricos.

Sean $x=A_{\rm rad}$ e $y=C^*_{\rm rad}$, con $x>|y|$, las coordenadas
levantadas ya construidas, y consérvense sus etiquetas de hoja. Escribimos
\begin{equation}
 q_\pm=e^{-(x\pm y)},\qquad s_\pm=q_\pm^{30},\qquad
 r_\pm=f(s_\pm),\qquad
 f(s)=\frac{1+s+s^2}{1+s+s^2+s^3}.
 \label{iii:eq:forma-canales}
\end{equation}
El dominio contractivo da $s_\pm\in(0,1)$ y
$r_\pm\in(3/4,1)$. En la cámara $x>y>0$ se cumplen
$s_+<s_-$ y $r_+>r_-$; el intercambio de hojas transporta esta cámara
a la orientación opuesta.

La imagen de las dos coordenadas constitutivas normalizadas es el dominio
\begin{equation}
 \mathcal U=\left\{(z,v)\in(0,\infty)^2:
   (zv)^{-1/2}\in(3/4,1),\quad
   (z/v)^{1/2}\in(3/4,1)\right\}.
 \label{iii:eq:forma-dominio}
\end{equation}
Aquí $(z,v)=(\widehat Z,\widehat c)$. Las restricciones expresan la
pertenencia de los canales reconstruidos al dominio del teorema de
inversión; no constituyen una elección ulterior de sus valores.

\Needspace{23\baselineskip}
\begin{proposition}[Composición constitutiva de la forma y de la impedancia LC]
\label{iii:prop:forma-LC}
Para $(\widehat Z,\widehat c)\in\mathcal U$, sean $s_\pm\in(0,1)$ las
raíces únicas de
\begin{equation}
 r_\pm s_\pm^3+(r_\pm-1)(1+s_\pm+s_\pm^2)=0,
 \qquad
 r_+=\frac1{\sqrt{\widehat Z\widehat c}},\quad
 r_-=\sqrt{\frac{\widehat Z}{\widehat c}}.
 \label{iii:eq:forma-inversion}
\end{equation}
La razón angular, la anisotropía de la elipse de acción y la impedancia
LC relativa se reconstruyen exactamente mediante
\begin{align}
 \frac{y}{x}
 &=\frac{\log s_+-\log s_-}{\log s_++\log s_-},
 \label{iii:eq:forma-razon}\\
 \eta_{\rm el}
 &=\frac12\log\frac{\log s_+}{\log s_-},
 \label{iii:eq:forma-eta}\\
 \frac{Z_{\rm LC}}{Z_*}
 &=e^{-\eta_{\rm el}}
  =\sqrt{\frac{\log s_-}{\log s_+}}.
 \label{iii:eq:forma-LC}
\end{align}
En la última identidad $Z_{\rm LC}=Z_{\eta_{\rm el}}$ es la impedancia
de la carta LC de \eqref{hmtII:eq:ii-elipse-lc}.
\end{proposition}
\begin{proof}
La inversión cúbica de \eqref{iii:eq:cubic} determina $s_\pm$ de manera
unívoca, porque $f$ es una biyección decreciente de $(0,1)$ sobre
$(3/4,1)$. Sus raíces positivas de orden $30$ recuperan $q_\pm$.
Las identidades
\[
 \log s_+=-30(x+y),\qquad \log s_-=-30(x-y)
\]
prueban \eqref{iii:eq:forma-razon} por diferencia y suma. Ambos logaritmos
son negativos; su cociente es positivo y su suma no es nula. Por ello,
\[
 \frac12\log\frac{\log s_+}{\log s_-}
   =\frac12\log\frac{x+y}{x-y}
   =\operatorname{artanh}(y/x)=\eta_{\rm el}.
\]
La relación $Z_{\eta_{\rm el}}/Z_*=e^{-\eta_{\rm el}}$, demostrada en
\eqref{hmtII:eq:ii-elipse-impedancia}, da \eqref{iii:eq:forma-LC}.
\end{proof}

La composición conserva la información necesaria para cada reconstrucción.
Con $\hbar$ se recuperan
$a_S=\sqrt{2\hbar}\,e^{\eta_{\rm el}/2}$ y
$b_S=\sqrt{2\hbar}\,e^{-\eta_{\rm el}/2}$; con la carta
$(\omega,Z_*)$ se recuperan la inductancia y la capacitancia ya definidas.
La acción, el reloj y la referencia dimensional acompañan la forma como
datos tipados. La inversión de dos coordenadas adimensionales no
reconstruye sus unidades ni toda la memoria del estado.

La normalización interviene antes de resolver las ecuaciones cúbicas.
Si se reciben $Z_{\rm vac}$ y $c_{\rm vac}$ en una carta dimensional,
las relaciones \eqref{iii:eq:dimensions} determinan primero
\[
 \widehat Z=\frac{Z_{\rm vac}}{u_Su_Q^{-2}},\qquad
 \widehat c=\frac{c_{\rm vac}}{u_C}.
\]
Esas son las coordenadas a las que se aplica
\eqref{iii:eq:forma-inversion}. La referencia $Z_*$ pertenece a la
realización LC y se conserva expresamente en \eqref{iii:eq:forma-LC}.

Los dos cocientes están relacionados mediante los canales, pero evalúan
funciones diferentes:
\begin{equation}
 \widehat Z=\frac{f(s_-)}{f(s_+)},\qquad
 \frac{Z_{\rm LC}}{Z_*}
       =\sqrt{\frac{\log s_-}{\log s_+}}.
 \label{iii:eq:forma-dos-cocientes}
\end{equation}
La reconstrucción anterior proporciona el mapa entre ambas lecturas.
En particular, intercambiar las hojas invierte ambos cocientes, cambia
$\eta_{\rm el}$ por $-\eta_{\rm el}$ y conserva $\widehat c$.
En el límite simétrico $y=0$, los canales coinciden y ambos cocientes
valen $1$. Fuera de ese caso, la igualdad de origen no impone igualdad
de las dos funciones. La forma de acción es recuperable desde la
respuesta constitutiva conservando orientación y normalización.
